\begingroup
\renewcommand{\arraystretch}{1.45}
\fancyhead[L]{\small LEMNISCATIC FUNCTIONS II.}
\newcommand{\GaussIIITranscriptionDigits}[1]{\text{\ttfamily #1}}
\SourceLeaf{435}

\par\medskip\Needspace{5\baselineskip}\pdfbookmark[0]{Collection of Calculations}{sammlung-rechnungen}\centerline{[17.]}

\begin{center}\textit{Collection of Calculations,}\end{center}
chiefly those employing my methods for finding the factors of large
numbers and \textsc{Wolfram's} logarithm tables.

\Needspace{26\baselineskip}
\begin{center}\textit{First Calculation for $e^{-\pi}=A$}\end{center}

A preliminary approximation had already shown that $A$ was, to the
$14^{\mathrm{th}}$ digit, $0,0432139182\ 6377$; hence $11.10^{10}A$ was very nearly $4753531008$
$=128.243.152827$. The question was therefore whether the number $152827$ could be resolved further into simple
factors.

\Needspace{26\baselineskip}
Division by small primes was unsuccessful; more elaborate methods therefore
had to be used. Comparing the number itself and various multiples of it with
the nearest squares gave, among other things,
\[
552^2\equiv950,\quad677^2\equiv152\quad\mathrm{mod.}\ 152827
\]
from which it was immediately concluded that
\[
1104^2\equiv3800,\quad3385^2\equiv3800
\]
and thus that the number $152827$ was not prime, since otherwise the squares of two
numbers both smaller than half of it could not possibly be congruent
(\textit{Disqu. Arr. art.}\quad.). The method taught in that work
(\textit{art.}\quad) then gave the factors of the number $152827$, namely $67.2281$.

\Needspace{26\baselineskip}
It was therefore certain that the hyperbolic logarithm of
\[
N=128.243.67.2281.11^{-1}.10^{-10}
\]
could be determined from the available tables and, at the same time, that it could
differ from the given logarithm, $-\pi$, only in the $10^{\mathrm{th}}$ decimal place.
It was therefore only necessary to calculate the number whose logarithm is this
difference and multiply it by $N$ to obtain $A$,
\[
A=Ne^{\log11+10\log10-\log2281-\log2144-\log972-\pi}=Ne^\delta
\]

\Needspace{34\baselineskip}
\SourceLeaf{436}

\textsc{Wolfram's} table gave:

[For brevity, set $\delta-\log(1+2135.10^{-13})=\delta'$,
$\delta'-\log(1+1443.10^{-17})=\delta''$.]

\begingroup\footnotesize\setlength{\arraycolsep}{1pt}\renewcommand{\arraystretch}{1.12}\[
\begin{array}{rcrlllll}
10\log10&&&&&&&\\[5pt]
&=&\GaussIIITranscriptionDigits{23,0258509299}&\GaussIIITranscriptionDigits{4045684017}&\GaussIIITranscriptionDigits{9914546843}&\GaussIIITranscriptionDigits{6420760110}&\GaussIIITranscriptionDigits{14886288}&\\[5pt]
\log11&=&\GaussIIITranscriptionDigits{2,3978952727}&\GaussIIITranscriptionDigits{9837054406}&\GaussIIITranscriptionDigits{1943577965}&\GaussIIITranscriptionDigits{1292998217}&\GaussIIITranscriptionDigits{06853937}&\\[5pt]
&&\GaussIIITranscriptionDigits{25,4237462027}&\GaussIIITranscriptionDigits{3882738424}&\GaussIIITranscriptionDigits{1858124808}&\GaussIIITranscriptionDigits{7713758327}&\GaussIIITranscriptionDigits{21740225}&\\[5pt]
l.2281&=&\GaussIIITranscriptionDigits{7,7323692222}&\GaussIIITranscriptionDigits{8438803081}&\GaussIIITranscriptionDigits{0466064812}&\GaussIIITranscriptionDigits{2168095619}&\GaussIIITranscriptionDigits{53159812}&\\[5pt]
l.2144&=&\GaussIIITranscriptionDigits{7,6704285221}&\GaussIIITranscriptionDigits{9069260675}&\GaussIIITranscriptionDigits{6232603654}&\GaussIIITranscriptionDigits{6055909444}&\GaussIIITranscriptionDigits{33575023}&\\[5pt]
l.972&=&\GaussIIITranscriptionDigits{6,8793558044}&\GaussIIITranscriptionDigits{6043907581}&\GaussIIITranscriptionDigits{0690427528}&\GaussIIITranscriptionDigits{9816593884}&\GaussIIITranscriptionDigits{53057834}&\\[5pt]
\pi&=&\GaussIIITranscriptionDigits{3,1415926535}&\GaussIIITranscriptionDigits{8979323846}&\GaussIIITranscriptionDigits{2643383279}&\GaussIIITranscriptionDigits{5028841971}&\GaussIIITranscriptionDigits{69399375}&\\[5pt]
&&\GaussIIITranscriptionDigits{25,4237462025}&\GaussIIITranscriptionDigits{2531295184}&\GaussIIITranscriptionDigits{0032479275}&\GaussIIITranscriptionDigits{3069440920}&\GaussIIITranscriptionDigits{09192044}&\\[5pt]
\delta&=&\GaussIIITranscriptionDigits{0,.........2}&\GaussIIITranscriptionDigits{1351443240}&\GaussIIITranscriptionDigits{1825645533}&\GaussIIITranscriptionDigits{4644317407}&\GaussIIITranscriptionDigits{12548181}&\\[5pt]
\delta'&=&\GaussIIITranscriptionDigits{0,..........}&\GaussIIITranscriptionDigits{...1443242}&\GaussIIITranscriptionDigits{4616770530}&\GaussIIITranscriptionDigits{2204949495}&\GaussIIITranscriptionDigits{65316893}&\\[5pt]
\delta''&=&\GaussIIITranscriptionDigits{0,..........}&\GaussIIITranscriptionDigits{.......242}&\GaussIIITranscriptionDigits{4616770634}&\GaussIIITranscriptionDigits{3329449495}&\GaussIIITranscriptionDigits{6431533124}&\\[5pt]
\frac12\delta''\delta''&=&\GaussIIITranscriptionDigits{0,..........}&\GaussIIITranscriptionDigits{..........}&\GaussIIITranscriptionDigits{..........}&\GaussIIITranscriptionDigits{.....29393}&\GaussIIITranscriptionDigits{8324222063}&\\[5pt]
e^{\delta''}&=&\GaussIIITranscriptionDigits{1,..........}&\GaussIIITranscriptionDigits{.......242}&\GaussIIITranscriptionDigits{4616770634}&\GaussIIITranscriptionDigits{3329478889}&\GaussIIITranscriptionDigits{4755755187}&\\[5pt]
e^{\delta'}&=&\GaussIIITranscriptionDigits{1,..........}&\GaussIIITranscriptionDigits{...1443242}&\GaussIIITranscriptionDigits{4616770634}&\GaussIIITranscriptionDigits{3679351089}&\GaussIIITranscriptionDigits{4781097632}&\\[5pt]
e^\delta&=&\GaussIIITranscriptionDigits{1,.........2}&\GaussIIITranscriptionDigits{1351443242}&\GaussIIITranscriptionDigits{4619851957}&\GaussIIITranscriptionDigits{0236156393}&\GaussIIITranscriptionDigits{8536512211}&\\[5pt]
11A&=&\GaussIIITranscriptionDigits{0,4753531009}&\GaussIIITranscriptionDigits{0149474751}&\GaussIIITranscriptionDigits{8595108889}&\GaussIIITranscriptionDigits{0081240330}&\GaussIIITranscriptionDigits{0920791693}&\GaussIIITranscriptionDigits{5}\\[5pt]
A&=&\GaussIIITranscriptionDigits{0,0432139182}&\GaussIIITranscriptionDigits{6377224977}&\GaussIIITranscriptionDigits{4417737171}&\GaussIIITranscriptionDigits{7280112757}&\GaussIIITranscriptionDigits{2810981063}&
\end{array}
\]\endgroup

\Needspace{20\baselineskip}
To make sure of the correctness of this result by a second calculation,
the number $A$ was multiplied by $599.10^8$, which gave
$2588513703,999957..$, so that a very easy calculation remained
if the number $2588513704$ could be resolved into factors smaller than $10000$.
A trial gave $2588513704=8.7.17.2719027$. The question was
therefore whether $2719027$ was prime. It was found that $-1848$ must certainly
be a quadratic residue of $2719027$ if this number was prime,
and that in that case this number must occur \textit{once} in the form $3xx+616yy$ and, conversely, that if composite it must be representable by this form
either not at all or \textit{more than once}.
But the method of exclusion showed that the number did indeed occur only once
in the form $3xx+616yy$, namely $2719027=3.197^2+616.65^2$,
\SourceLeaf{437}
from which it followed with certainty that $2719027$ was prime and that the
method attempted was therefore not applicable this time.

\Needspace{18\baselineskip}
After a few other unsuccessful attempts, the following method was finally
tried, with better success. The number $A$ multiplied by $10^{14}$ gave $4321391826377,25$
and the number $4321391826375$ resolved into the factors $125.81.13.32831087$. That
the number $32831087$ was not prime followed from its not occurring in the
form $xx+190yy$; the first of the methods taught in the \textit{Disqui. Arr.}
was therefore applied to it, revealing that it was the product
of $373.8819$.

\Needspace{30\baselineskip}
\begin{center}\textit{Second Calculation for $e^{-\pi}=A$}\end{center}

$59.10^9A$ was found very nearly $=2549621178=54.13.3631939$. The number
$3631939$ resolved into the factors $1091.3329$ (which had been found from
\[
3631939=40.232^2+19.279^2=40.300^2+19.41^2
\]
). Thus
\[
\begin{aligned}59A&=702.1091.3329.10^{-9}\\[7pt]&\quad{}e^{9\log10+\log59-\log702-\log1091-\log3329-\pi}=Ne^\delta\end{aligned}
\]
\Needspace{29\baselineskip}
\textsc{Wolfram's} table gave [setting $\delta-\log(1-17157.10^{-14})=\delta'$,
\hfill\break{} $\frac12\delta'\delta'+\frac16\delta'^3=\varepsilon$]
\begingroup\footnotesize\setlength{\arraycolsep}{1pt}\renewcommand{\arraystretch}{1.12}\[
\begin{array}{rcrlllll}
C9l.10&=&\GaussIIITranscriptionDigits{21,2767341630}&\GaussIIITranscriptionDigits{5358884383}&\GaussIIITranscriptionDigits{8076907840}&\GaussIIITranscriptionDigits{7221315900}&\GaussIIITranscriptionDigits{86602341}&\\[5pt]
C.l.59&=&\GaussIIITranscriptionDigits{5,9224625560}&\GaussIIITranscriptionDigits{9428054938}&\GaussIIITranscriptionDigits{3949626280}&\GaussIIITranscriptionDigits{3023759366}&\GaussIIITranscriptionDigits{53210670}&\\[5pt]
l.702&=&\GaussIIITranscriptionDigits{6,5539334040}&\GaussIIITranscriptionDigits{2581111965}&\GaussIIITranscriptionDigits{6455273791}&\GaussIIITranscriptionDigits{0722868232}&\GaussIIITranscriptionDigits{39752589}&\\[5pt]
l.1091&=&\GaussIIITranscriptionDigits{6,9948499858}&\GaussIIITranscriptionDigits{3307081851}&\GaussIIITranscriptionDigits{1895817110}&\GaussIIITranscriptionDigits{3840806982}&\GaussIIITranscriptionDigits{08318537}&\\[5pt]
l.3329&=&\GaussIIITranscriptionDigits{8,1104272375}&\GaussIIITranscriptionDigits{7502494295}&\GaussIIITranscriptionDigits{2021653586}&\GaussIIITranscriptionDigits{1576658324}&\GaussIIITranscriptionDigits{67001597}&\\[5pt]
\pi&=&\GaussIIITranscriptionDigits{3,1415926535}&\GaussIIITranscriptionDigits{8979323846}&\GaussIIITranscriptionDigits{2643383279}&\GaussIIITranscriptionDigits{5028841971}&\GaussIIITranscriptionDigits{69399375}&\\[5pt]
-\delta&=&\GaussIIITranscriptionDigits{0,.........1}&\GaussIIITranscriptionDigits{7156951280}&\GaussIIITranscriptionDigits{5042661888}&\GaussIIITranscriptionDigits{1414250778}&\GaussIIITranscriptionDigits{24285109}&\\[5pt]
\delta'&=&\GaussIIITranscriptionDigits{0,..........}&\GaussIIITranscriptionDigits{.....48720}&\GaussIIITranscriptionDigits{9675470563}&\GaussIIITranscriptionDigits{5420349120}&\GaussIIITranscriptionDigits{2333831382}&\\[5pt]
\varepsilon&=&\GaussIIITranscriptionDigits{0,..........}&\GaussIIITranscriptionDigits{..........}&\GaussIIITranscriptionDigits{..........}&\GaussIIITranscriptionDigits{1186866339}&\GaussIIITranscriptionDigits{3606594227}&\\[5pt]
e^{\delta'}&=&\GaussIIITranscriptionDigits{1,..........}&\GaussIIITranscriptionDigits{.....48720}&\GaussIIITranscriptionDigits{9675470563}&\GaussIIITranscriptionDigits{6607215459}&\GaussIIITranscriptionDigits{5940425609}&\\[5pt]
1-e^\delta&=&\GaussIIITranscriptionDigits{0,.........1}&\GaussIIITranscriptionDigits{7156951279}&\GaussIIITranscriptionDigits{0324613026}&\GaussIIITranscriptionDigits{9032989386}&\GaussIIITranscriptionDigits{4786573955}&\\[5pt]
59A&=&\GaussIIITranscriptionDigits{2,5496211775}&\GaussIIITranscriptionDigits{6256273669}&\GaussIIITranscriptionDigits{0646493131}&\GaussIIITranscriptionDigits{9526652679}&\GaussIIITranscriptionDigits{5847882752}&\GaussIIITranscriptionDigits{6}\\[5pt]
A&=&\GaussIIITranscriptionDigits{0,0432139182}&\GaussIIITranscriptionDigits{6377224977}&\GaussIIITranscriptionDigits{4417737171}&\GaussIIITranscriptionDigits{7280112757}&\GaussIIITranscriptionDigits{2810981063}&\GaussIIITranscriptionDigits{6}
\end{array}
\]\endgroup

\Needspace{9\baselineskip}
\SourceLeaf{438}

At the same time, the agreement of both calculations allows the conclusion
that the logarithms of $2,3,5,11,13,59,67,1091,2281,3329$ in
\textsc{Wolfram's} table are correct to the last digit.

\clearpage
\endgroup
